\chapter*{كلمة المحرر}
\addcontentsline{toc}{chapter}{كلمة المحرر}

الحمد لله الذي بنعمته تتم الصالحات، والصلاة والسلام على المبعوث رحمة للعالمين، نبينا محمد وعلى آله وصحبه أجمعين، أما بعد؛

فإن كتاب «الرسالة» للإمام الجليل أبي عبد الله محمد بن إدريس الشافعي (المتوفى سنة 204 هـ) يمثل قطب الرحى في علوم الشريعة، واللبنة الأولى التي أُسس عليها علم أصول الفقه، ولا يستغني عنه طالب علم ولا فقيه. وقد يسر الله تعالى لفضيلة الشيخ سامي -حفظه الله- شرح هذا السفر العظيم وتوضيح مقاصده، وفك عباراته في مجالس علمية عامرة بالفوائد والدرر.

وحرصاً منا على تعميم النفع بهذه المجالس، وحفظاً لهذا الشرح النفيس من الضياع، قمنا بالاعتناء بتفريغات هذه الدروس وإخراجها في حلة مطبوعة تليق بمكانة الكتاب وشرحه، وقد اتبعنا في إعداد هذا الكتاب المنهجية التالية:

\begin{enumerate}
    \item \textbf{التنقيح اللغوي وإعادة الصياغة:} قمنا بتحويل لغة الشرح الشفهية المفرغة إلى لغة عربية فصحى مقروءة وسلسة، مع إزالة اللزمات الكلامية، وحذف الاستطرادات العابرة التي تخرج عن مقصود الشرح، وذلك حفاظاً على التسلسل المنطقي للأفكار، مع الحرص التام على عدم الإخلال بالمعنى العلمي الذي قصده الشيخ.
    \item \textbf{التنسيق والتبويب:} تم تقسيم مادة الكتاب إلى فصول ومباحث ذات عناوين واضحة ودقيقة تعكس مضمون كل قسم، مما يسهل على القارئ فهم الهيكلة العامة والعودة إلى المسائل بيسر.
    \item \textbf{العناية بالنصوص الشرعية:} تم تمييز الآيات القرآنية بالأقواس المخصصة لها، وكتابتها بالرسم المناسب، وتنسيق الأحاديث النبوية ووضعها بين علامات التنصيص لتبرز عن كلام الشارح.
    \item \textbf{استخلاص الفوائد والملاحق:} لم نكتفِ بمجرد تنسيق الشرح، بل قمنا بعد كل درس باستخلاص أهم الفوائد والدرر التي نثرها الشيخ طيات شرحه، وقمنا بتبويبها وتصنيفها موضوعياً (عقدياً، ومنهجياً، وأصولياً، ولغوياً) وإدراجها في ملاحق خاصة في نهاية الكتاب، لتكون زاداً سريعاً للقارئ ومراجعة مركزة لمقاصد الشرح.
\end{enumerate}

وقد اعتمدنا في تنسيق الكتاب وإخراجه فنياً على نظام التنضيد الاحترافي (\LaTeX) لضمان أعلى معايير الجودة الطباعية.

وختاماً، نسأل الله العلي القدير أن يجعل هذا العمل خالصاً لوجهه الكريم، وأن يثيب فضيلة الشيخ على ما قدم، وأن ينفع بهذا الكتاب كل من قرأه أو ساهم في إخراجه، إنه سميع مجيب.

\vspace{1cm}
\noindent
\textbf{المحرر} \\
\textbf{فريق العمل}